\newpage
\subsection{p-Wert}
Für Stichprobe $\cX_i$ ($n$ el.) testen $H_0 : \vartheta = \vartheta_0$ gegen eine $H_A : \vartheta \in \Theta_A$.

$p$-Wert ist W. unter $H_0$ einen mind. so extremen Wert für $T$ zu bekommen wie $T(\omega) = t(x_1, \ldots, x_n)$ (extrem bezieht sich auf $H_A$).
Für bspw. $H_A: \vartheta > \vartheta_0$, dann ist $p\text{-Wert}(\omega) = \P_{\vartheta_0}[T > t_0] \Big|_{t_0 = t(\cX_1(\omega), \ldots)}$

\shortdefinition[Geordnete Testsammlung]
$T$ Teststatistik, Familie von Tests $(T, K_t)_{t \geq 0}$ \bi{geordnet bezgl. $T$}, falls $\forall t$ gilt $K_t \subseteq \R$ und $\forall s \leq t$ gilt $K_s \supset K_t$.
Typische Beispiele:
\begin{itemize}
    \item $K_t = (t, \8)$ (rechtsseitiger Test)
    \item $K_t = (-\8, t)$ (linksseitiger Test)
    \item $K_t = (-\8, t) \cup (t, \8)$ (beidseitiger Test)
\end{itemize}

\shortdefinition[p-Wert] $p\text{-Wert} = G(T) : \Omega \rightarrow [0, 1]$, mit $G(t) = \P_{\vartheta_0}[T \in K_t]$.
Er ist eine Zufallsvariable, hängt direkt von $x_i$ ab. Wiederholen des Tests generiert neuen $p$-Wert.

Wenn $T$ stetig und $K_t = (t, \8)$, dann gleichverteilt auf $[0, 1]$.

\inlineintuition Besagt, welche Tests $H_0$ ablehnen würden (alle Tests mit $\alpha > p$ lehnen $H_0$ ab). $H_A$ spielt \textit{keine} Rolle.
Es heisst \textit{NICHT}, dass der $p$-Wert die W. \textit{IST}, dass die Hypothese richtig ist.
